\documentclass[11pt]{article}
\usepackage[margin=1in]{geometry}
\usepackage{amsmath,amssymb,amsthm,booktabs,hyperref,listings}
\newtheorem{theorem}{Theorem}
\newtheorem{lemma}{Lemma}
\newcommand{\Q}{\mathbb Q}
\title{Four roots of Yuster's entropy polynomial at $k=50$}
\author{Principia Math In Progress}
\date{18 September 2026}
\begin{document}
\maketitle
\begin{abstract}
An exact rational sign certificate disproves Conjecture 3.9 of Yuster:
$p_{50}$ has at least four distinct roots in $(0,1)$. The certificate uses
integer arithmetic and a standard multisection formula, also appearing in
Wakhare's work. It concerns the polynomial root conjecture; the associated
entropy inequality has already been proved by Ho.
\end{abstract}

\section*{Task 1: complete disproof}
For an integer $k\ge2$, let $\varphi_k\in(0,1)$ satisfy
$\varphi_k^k+\varphi_k=1$, and put
$\alpha_k=\varphi_k^{k-1}=\varphi_k^{-1}-1$.
Equivalently, $\alpha_k$ is the unique positive solution of
\begin{equation}\label{alpha}
 \alpha_k(1+\alpha_k)^{k-1}=1.
\end{equation}
Set $h(x)=-x\log x-(1-x)\log(1-x)$ and
$f_k(x)=\alpha_k h(x^k)-x^{k-1}h(x)$ on $(0,1)$.
Yuster's polynomial is characterized by his Corollary 3.7 \cite{Y}:
\begin{equation}\label{source}
 f_k^{(k+1)}(x)=\frac{(k-1)!p_k(x)}{x(x^k-1)^k}.
\end{equation}
This is exactly the normalization in the original conjecture.

\begin{theorem}\label{main}
The polynomial $p_{50}$ has a root in each of
\[
 (0,1/4),\qquad(1/4,1/2),\qquad(1/2,3/5),\qquad(3/5,1).
\]
In particular, the assertion that every $p_k$ has at most two roots in
$(0,1)$, counting multiplicity, is false.
\end{theorem}

\subsection*{A convenient exact polynomial formula}
Define
\[
 S_k(x)=\sum_{i=0}^{k-1}x^i,\quad
 B_k(x)=\sum_{i=0}^{k-1}(1-x)^i,\quad
 H_k(y)=\sum_{j=0}^{k-1}a_{k,j}y^j,
\]
where
\begin{equation}\label{coef}
 a_{k,j}=[z^{kj}]S_k(z)^k
 =\sum_{v=0}^{j}(-1)^v\binom{k}{v}
          \binom{k(j-v)+k-1}{k-1}.
\end{equation}
Every $a_{k,j}$ is positive: $kj$ can be written as a sum of $k$ integers
between $0$ and $k-1$. In particular $H_k(y)>0$ for $y\ge0$.
These coefficients and the following derivative identities are known;
see Wakhare's Corollaries 13 and 14 \cite{W}. We include the derivation
to fix all signs and normalizing factors.

\begin{lemma}\label{formula}
One has the polynomial identity
\begin{equation}\label{p}
 p_k(x)=(-1)^k\bigl(B_k(x)S_k(x)^k-\alpha_k k^2H_k(x^k)\bigr).
\end{equation}
\end{lemma}
\begin{proof}
Taking coefficients at multiples of $k$ in
$(1-z)^{-k}=S_k(z)^k/(1-z^k)^k$ gives
\begin{equation}\label{section}
 \sum_{n\ge0}\binom{kn+k-1}{k-1}x^{kn}
       =\frac{H_k(x^k)}{(1-x^k)^k}.
\end{equation}
The convergent expansion
$h(u)=-u\log u+u-\sum_{n\ge1}u^{n+1}/(n(n+1))$
may be differentiated term by term on compact subintervals of $(0,1)$.
It yields
\begin{align*}
 \frac{d^{k+1}}{dx^{k+1}}h(x^k)
   &=-\frac{k\,k!}{x}\sum_{n\ge0}\binom{kn+k-1}{k-1}x^{kn},\\
 \frac{d^{k+1}}{dx^{k+1}}\bigl(x^{k-1}h(x)\bigr)
   &=\frac{(k-1)!}{x^2}\bigl(1-(1-x)^{-k}\bigr).
\end{align*}
Combining these equations with \eqref{section}, using
$B_k(x)=(1-(1-x)^k)/x$ and
$S_k(x)=(1-x^k)/(1-x)$ on $(0,1)$, gives
\[
 f_k^{(k+1)}(x)=\frac{(k-1)!}{x(1-x^k)^k}
 \bigl(B_k(x)S_k(x)^k-\alpha_k k^2H_k(x^k)\bigr).
\]
Comparison with \eqref{source} proves \eqref{p} on $(0,1)$, hence
identically as a polynomial.
\end{proof}

\subsection*{Exact signs}
For $0\le x\le1$, put
\[
 R_k(x)=\frac{B_k(x)S_k(x)^k}{k^2H_k(x^k)}.
\]
The denominator is positive. For even $k$, the sign of $p_k(x)$ is
the sign of $R_k(x)-\alpha_k$.
Since $t(1+t)^{49}$ is strictly increasing for $t>0$, the integer
inequalities
\[
 59\cdot1059^{49}<1000^{50},\qquad
 3\cdot53^{49}>50^{50}
\]
give
\begin{equation}\label{bracket}
 \frac{59}{1000}<\alpha_{50}<\frac{3}{50}.
\end{equation}
The following rational bounds follow directly from \eqref{coef}:
\begin{center}
\begin{tabular}{ccc}
\toprule
$x$ & Exact value or strict rational enclosure of $R_{50}(x)$
 & $\operatorname{sgn}p_{50}(x)$\\
\midrule
$0$ & $1/50$ & $-$\\
$1/4$ & $(2717,2718)$ & $+$\\
$1/2$ & $(1/50,21/1000)$ & $-$\\
$3/5$ & $(31/500,63/1000)$ & $+$\\
$1$ & $1/50$ & $-$\\
\bottomrule
\end{tabular}
\end{center}
For completeness, the entire arithmetic certificate for this table is
reproducible by the following Python code, using exact fractions only.
Thus none of the strict signs depends on a numerical root finder.
\begin{lstlisting}[basicstyle=\ttfamily\small,frame=single]
from fractions import Fraction as Q
from math import comb
k = 50
a = [sum((-1)**v*comb(k,v)*comb(k*(j-v)+k-1,k-1)
         for v in range(j+1)) for j in range(k)]
def R(x):
    S = sum(x**i for i in range(k))
    B = sum((1-x)**i for i in range(k))
    H = sum(Q(c)*x**(k*j) for j,c in enumerate(a))
    return B*S**k/(k*k*H)
assert 59*1059**49 < 1000**50
assert 3*53**49 > 50**50
assert R(Q(0)) == R(Q(1)) == Q(1,50)
assert 2717 < R(Q(1,4)) < 2718
assert Q(1,50) < R(Q(1,2)) < Q(21,1000)
assert Q(31,500) < R(Q(3,5)) < Q(63,1000)
\end{lstlisting}
The accompanying \texttt{exact-certificate.json} stores the coefficients
and exact fractions. The separate \texttt{verify.py} additionally
reconstructs the coefficients by bounded-composition convolution and
checks the original $C$ recurrence and $\rho_k,\sigma_k$ sums from
\cite{Y} at all five points, including $k=50$.

\begin{proof}[Proof of Theorem~\ref{main}]
The table and \eqref{bracket} give the strict sign sequence
$-,+,-,+,-$. The intermediate value theorem, applied on each of the
four adjacent closed intervals, gives a root in its interior. The four
open intervals are disjoint. These are therefore four distinct roots
of the polynomial specified by \eqref{source}, proving the assertion.
\end{proof}

\subsection*{Scope and prior results}
Yuster proved the root claim for $k=2,3,4$ and used it as a sufficient
condition for an entropy inequality. Ho \cite{H} later proved that
entropy inequality for every real $k>1$ by a different argument.
The present counterexample concerns the sufficient root condition.
The derivative and multisection identities used above are explicitly
credited to the prior literature. No claim is made here that $50$ is
the least counterexample or that $p_{50}$ has exactly four roots.
Wakhare's polynomial with parameters $(k,r)=(50,1)$ differs from
$p_{50}$ by the nonzero factor $-(1-x)^{50}/50$ on $(0,1)$; thus the
same witness also refutes his broader Conjecture 2. This consequence
receives no separate campaign credit.

\section*{Tasks 2--5: proposed follow-up work}
\begin{enumerate}
\setcounter{enumi}{1}
\item Determine the least integer $k\ge2$ for which $p_k$ has more than
two roots in $(0,1)$, with an exact certificate for every smaller $k$.
\item Prove that the number of distinct roots of $p_k$ in $(0,1)$ is
unbounded as $k\to\infty$.
\item Find asymptotic locations of the successive root pairs, using the
multisection sum in \eqref{section}.
\item Classify the coprime integer pairs $k>r\ge1$ for which Wakhare's
entropy polynomial $p_{k,r}$ has exactly two roots in $(0,1)$.
\end{enumerate}

\begin{thebibliography}{9}
\bibitem{Y} R. Yuster, \emph{Almost $k$-union closed set systems},
arXiv:2302.12276v1 (2023), especially Corollary 3.7 and Conjecture 3.9.
\url{https://arxiv.org/abs/2302.12276}.
\bibitem{W} T. Wakhare, \emph{Iterated entropy derivatives and binary
entropy inequalities}, Journal of Approximation Theory \textbf{307}
(2025), 106143. Complete arXiv version read:
\url{https://arxiv.org/abs/2312.14743}, v2, Corollaries 13--14.
\bibitem{H} B. S. Ho, \emph{A generalization of Boppana's entropy
inequality}, arXiv:2601.19327v1 (2026).
\url{https://arxiv.org/abs/2601.19327}.
\end{thebibliography}
\end{document}
